\section{Equipo y contribución individual}\label{sec:equipo}

\kyrian{} es el trabajo de cuatro personas que se repartieron el proyecto por bloques: cada
miembro lideró al menos uno de principio a fin y colaboró en los de los demás. La
Tabla~\ref{tab:raci} detalla quién fue responsable de cada componente y quién contribuyó, agrupando
los componentes por área, y la Fig.~\ref{fig:fases} muestra cómo se fue repartiendo el trabajo a lo
largo del proyecto. Todo el equipo se apoyó en asistentes de programación (sobre todo Claude Code,
el asistente de programación de Anthropic); la autoría que se recoge aquí refleja quién especificó,
dirigió, revisó e integró cada parte.

\begin{table}[!h]
\centering\scriptsize
\caption{Matriz de responsabilidad por componente, agrupada por área (R: responsable; C: contribuye). MP: Manuel Pérez; JC: Joaquín Castro Salas; DB: David Baos; AD: Alexander De Sousa.}\label{tab:raci}
\begin{tabularx}{\columnwidth}{@{}Ycccc@{}}
\toprule
\textbf{Componente} & \textbf{MP} & \textbf{JC} & \textbf{DB} & \textbf{AD} \\
\midrule
\multicolumn{5}{@{}l}{\textit{Producto y negocio}} \\
Idea del producto y estudio de mercado & C & & & R \\
Planificador web: interfaz, mapa y administración & R & C & & C \\
Diseño visual en torno a Kiri & & R & & \\
\midrule
\multicolumn{5}{@{}l}{\textit{Inteligencia artificial}} \\
Agente Kiri: herramientas, \emph{guardrails} y salidas estructuradas & R & & & C \\
Recuperación (RAG) con Bedrock y S3 Vectors & C & R & & \\
Ingesta de datos: corpus, puerta de calidad y ciudades & R & C & C & C \\
Evaluación: recuperación, respuestas y agente & R & R & & \\
\midrule
\multicolumn{5}{@{}l}{\textit{Arquitectura y nube}} \\
Arquitectura de software y repositorio & R & & & \\
Servicio de datos: usuarios, viajes, acceso y cuota & R & C & & \\
Infraestructura en AWS con Terraform & R & & C & C \\
Borde HTTPS del frontend & C & & & R \\
\midrule
\multicolumn{5}{@{}l}{\textit{DevOps y operación}} \\
CI/CD y entorno de desarrollo & R & C & & C \\
Observabilidad: trazas, registros y cuadros de mando & R & & C & \\
Proceso de entrega con agentes & R & & & \\
\bottomrule
\end{tabularx}
\end{table}

\begin{figure}[t]
\centering
\resizebox{\linewidth}{!}{\begin{tikzpicture}[x=1cm,y=1cm,
  ph/.style={box,font=\scriptsize,text width=2.45cm,minimum height=0.55cm,fill=klight},
  tx/.style={font=\scriptsize,text width=2.5cm,align=left,anchor=north}]
\draw[arr,line width=1pt] (-1.4,0) -- (7.0,0);
\node[ph] (p1) at (0,0) {\textbf{Exploración}\\marzo--6 sep.};
\node[ph] (p2) at (2.8,0) {\textbf{Nueva arquitectura}\\7--20 sep.};
\node[ph] (p3) at (5.6,0) {\textbf{Producto}\\21 sep.--octubre};
\node[tx] at (0,-0.5) {AD: idea, estudio de mercado, primera fuente de datos, \emph{devcontainer} e IaC\\MP: web y frontend\\JC: backend inicial y primer \emph{chatbot}};
\node[tx] at (2.8,-0.5) {MP: dos servicios, \emph{serverless}, corpus y planificador\\JC: Bedrock y recuperación RAG\\AD: HTTPS del frontend};
\node[tx] at (5.6,-0.5) {MP: datos, trazas, acceso y agente\\JC: diseño visual y evaluación\\AD: salidas estructuradas\\DB: cuadros de mando};
\end{tikzpicture}}
\caption{Fases del proyecto y bloques liderados por cada miembro. Las dos decisiones que más cambiaron el sistema fueron pasar de contenedores a \emph{serverless} y de un planificador por intenciones a un agente con herramientas (§\ref{sec:evalagente}).}\label{fig:fases}
\end{figure}

Cada miembro lideró un bloque del proyecto:
\begin{itemize}
  \item \textbf{Alexander De Sousa}: la idea de \kyrian{} y el estudio de mercado que la sitúa
        frente a la competencia (§\ref{sec:negocio}); la primera fuente de datos, el primer
        entorno de desarrollo en contenedor y la primera infraestructura como código; el borde HTTPS
        del frontend (S3, CloudFront, Route~53 y ACM), y las salidas estructuradas nativas que usan
        el agente y los pasos deterministas del planificador.
  \item \textbf{Manuel Pérez}: la arquitectura de software y de nube y su infraestructura como
        código, el agente Kiri con sus herramientas y \emph{guardrails}, la ingesta de datos con su
        puerta de calidad y \code{/add-city}, la evaluación del agente, la observabilidad de la IA y
        el proceso de entrega con agentes.
  \item \textbf{Joaquín Castro Salas}: la recuperación aumentada (integración con Bedrock y
        S3 Vectors), la evaluación de la recuperación y de las respuestas con un juez LLM (un
        modelo de lenguaje que puntúa las respuestas) y el diseño visual de la interfaz en torno a
        Kiri.
  \item \textbf{David Baos}: los cuadros de mando de CloudWatch y las métricas de negocio sobre los
        registros de producción, y la configuración de una de las ciudades del corpus.
\end{itemize}
